\section*{Capítulo \ref{ch:b1:gene-expression} --- Expressão gênica: transcrição e tradução}

\begin{solution}{exo:b1:gene-expression:1}
\omterm{def:b1:nucleic-acids:chain}{RNA} $5'$-AUGCCUAAG-$3'$; fita codificante $5'$-ATGCCTAAG-$3'$.
\end{solution}

\begin{solution}{exo:b1:gene-expression:2}
Um capuz $5'$ (guanina modificada), uma \omterm{prop:b1:gene-expression:processing}{cauda poli-A} $3'$, e a remoção
dos \omterm{def:b1:genomes:gene}{íntrons} por splicing.
\end{solution}

\begin{solution}{exo:b1:gene-expression:3}
AUG CCG AAA GUU UGA: Met-Pro-Lys-Val, e então parada.
\end{solution}

\begin{solution}{exo:b1:gene-expression:4}
A: o tRNA carregado que chega é admitido e conferido. P: o tRNA que
carrega a cadeia em crescimento; a ligação se forma entre a cadeia dele
e o \omterm{def:b1:proteins:aminoacid}{aminoácido} do sítio A. E: o tRNA esvaziado a caminho da saída.
\end{solution}

\begin{solution}{exo:b1:gene-expression:5}
Codificantes $2400 - 600 = 1800$ \omterm{def:b1:nucleic-acids:nucleotide}{nucleotídeos}: 599 resíduos (600 \omterm{def:b1:gene-expression:code}{códons},
um deles de parada). Tempo $599/4 = \qty{150}{s}$. \omterm{def:b1:gene-expression:ribosome}{Ribossomos}
$1800/90 = 20$ ao mesmo tempo.
\end{solution}

\begin{solution}{exo:b1:gene-expression:6}
A mensagem é lida de três em três a partir de um início fixo; uma ou
duas bases a mais deslocam todos os \omterm{def:b1:gene-expression:code}{códons} seguintes e o resto da
\omterm{def:b1:proteins:peptide}{proteína} vira algaravia, em geral terminando numa parada prematura. Três
bases a mais acrescentam um \omterm{def:b1:gene-expression:code}{códon} e restauram o quadro: a \omterm{def:b1:proteins:peptide}{proteína} tem
um resíduo a mais e alguns errados entre as inserções, e muitas vezes
ainda funciona.
\end{solution}

\begin{solution}{exo:b1:gene-expression:7}
CAG (Gln) para UAG (parada): a cadeia termina no resíduo 150, uma
\omterm{def:b1:proteins:peptide}{proteína} truncada e quase certamente inativa (mutação sem sentido). CAG
para CAA: continua Gln, silenciosa. CAG para CGG: Arg no lugar de Gln,
uma substituição (mutação de sentido trocado) cujo efeito depende da
posição.
\end{solution}

\begin{solution}{exo:b1:gene-expression:8}
O \omterm{def:b1:gene-expression:ribosome}{ribossomo} confere apenas o pareamento entre \omterm{def:b1:gene-expression:code}{códon} e \omterm{def:b1:gene-expression:trna}{anticódon}; ele não
consegue ver o \omterm{def:b1:proteins:aminoacid}{aminoácido}. Se a cisteína presa ao tRNA dela for
convertida quimicamente em alanina, o \omterm{def:b1:gene-expression:ribosome}{ribossomo} insere alanina onde
quer que a mensagem diga cisteína: o que é lido é o tRNA, e não o
\omterm{def:b1:proteins:aminoacid}{aminoácido}. Logo, a sintetase, que emparelha cada \omterm{def:b1:proteins:aminoacid}{aminoácido} com o tRNA
certo, é a verdadeira tradutora.
\end{solution}

\begin{solution}{exo:b1:gene-expression:9}
$400\times 4 = 1600$ ligações; os dois GTP do \omterm{def:b1:gene-expression:ribosome}{ribossomo} por resíduo são
metade delas, e os dois das sintetases, a outra metade.
\end{solution}

\begin{solution}{exo:b1:gene-expression:10}
Nos eucariotos o transcrito é feito no núcleo e os \omterm{def:b1:gene-expression:ribosome}{ribossomos} ficam no
\omterm{def:b1:eukaryotic-cell:organelle}{citosol}, separados pelo envoltório; e o transcrito só é um mensageiro
depois de ter sofrido splicing e recebido o capuz. A separação é o que
torna possível o próprio processamento --- o splicing alternativo, o
controle da exportação, e uma verificação de que a mensagem está
completa antes de ser lida.
\end{solution}

\begin{solution}{exo:b1:gene-expression:11}
Quatro mensageiros: com os dois \omterm{def:b1:genomes:gene}{éxons} (190 \omterm{def:b1:nucleic-acids:nucleotide}{nucleotídeos} acrescentados),
só com o 3 (90), só com o 4 (100), sem nenhum dos dois (0). O quadro se
mantém se o comprimento acrescentado for múltiplo de 3: 0 e 90 o mantêm;
100 e 190 o deslocam, de modo que os mensageiros só com o \omterm{def:b1:genomes:gene}{éxon} 4 ou com
os dois \omterm{def:b1:genomes:gene}{éxons} leem o \omterm{def:b1:genomes:gene}{éxon} 5 fora do quadro e truncam a \omterm{def:b1:proteins:peptide}{proteína}. O
splicing alternativo precisa respeitar o quadro, e é por isso que os
\omterm{def:b1:genomes:gene}{éxons} muitas vezes são múltiplos de três.
\end{solution}

\begin{solution}{exo:b1:gene-expression:12}
Arbitrário: nada na química diz que UUU precise significar Phe; as
atribuições são convenções mantidas pelas sintetases. Otimizado: a
terceira posição é a mais degenerada, de modo que os erros e as mutações
que caem ali são muitas vezes silenciosos, e os \omterm{def:b1:gene-expression:code}{códons} que diferem por
uma base tendem a codificar \omterm{def:b1:proteins:aminoacid}{aminoácidos} parecidos, de modo que uma
substituição costuma ser branda --- o código minimiza o dano do erro.
Congelado: uma sintetase que mudasse a especificidade dela alteraria de
uma vez todas as \omterm{def:b1:proteins:peptide}{proteínas} da \omterm{def:b1:cell-unit-of-life:cell}{célula}, aos milhares, no mesmo instante;
quase nenhuma \omterm{def:b1:cell-unit-of-life:cell}{célula} sobrevive a isso, e por isso o código não pode
derivar e está fixado desde o ancestral comum de todos os seres vivos.
\end{solution}

\begin{solution}{pb:b1:gene-expression:1}
\textbf{1.} $30\,000/30 = \qty{1000}{s}$, cerca de \qty{17}{min}.
\textbf{2.} $28\,000/30\,000 = \qty{93}{\%}$.
\textbf{3.} $30\,000/2000 = 15$.
\textbf{4.} $2\times 30\,000 = \num{60000}$ \omterm{def:b1:water-small-molecules:atp}{ATP}.
\textbf{5.} $1000/10 = 100$ polimerases sobre o \omterm{def:b1:genomes:gene}{gene}; 360 mensageiros
por hora.
\textbf{6.} Sete \omterm{def:b1:genomes:gene}{íntrons}, com média de $28\,000/7 = \qty{4000}{}$
\omterm{def:b1:nucleic-acids:nucleotide}{nucleotídeos}.
\textbf{7.} A transcrição (\qty{17}{min}); os \omterm{def:b1:genomes:gene}{íntrons} sofrem splicing à
medida que são feitos, e o último acrescenta só um minuto.
\textbf{8.} \omterm{def:b1:genomes:gene}{Gene} $30\,000\times 0.34\,\mathrm{nm} = \qty{10}{\micro m}$; mensageiro \qty{0.68}{\micro m}.
\textbf{9.} $500/4 = \qty{125}{s}$.
\textbf{10.} $2000/80 = 25$.
\textbf{11.} Uma cadeia termina a cada $80/4 = \qty{20}{s}$: 180 por
hora.
\textbf{12.} $2/\ln 2 = \qty{2.9}{h}$ de produção plena: cerca de 520
\omterm{def:b1:proteins:peptide}{proteínas}.
\textbf{13.} $4\times 500 = 2000$ \omterm{def:b1:water-small-molecules:atp}{ATP}.
\textbf{14.} $60\,000/520 = 115$ \omterm{def:b1:water-small-molecules:atp}{ATP} por \omterm{def:b1:proteins:peptide}{proteína}.
\textbf{15.} Cerca de 2100 \omterm{def:b1:water-small-molecules:atp}{ATP}, \qty{95}{\%} deles na tradução.
\textbf{16.} $1 - (1 - 10^{-5})^{1503} \approx 1503\times 10^{-5} =
\qty{1.5}{\%}$.
\textbf{17.} $1 - (1 - 10^{-4})^{500} \approx \qty{5}{\%}$.
\textbf{18.} Transcrição: $1.5\%\times 0.75\times 0.67 \approx
0.75\%$ dos mensageiros defeituosos, e portanto das \omterm{def:b1:proteins:peptide}{proteínas};
tradução: $5\%\times 0.67 = 3.3\%$; cerca de \qty{4}{\%} das moléculas.
\textbf{19.} Um erro de \omterm{def:b1:proteins:peptide}{proteína} fica confinado a uma molécula, que logo
é degradada; um erro de mensageiro é copiado em centenas de \omterm{def:b1:proteins:peptide}{proteínas};
um erro de replicação é herdado por todos os descendentes para sempre. O
custo de um erro, e portanto a exatidão que vale a pena pagar, sobe a
cada passo para trás.
\textbf{20.} $3000/180 = 17$ mensageiros presentes.
\textbf{21.} $17\times 25 = 420$ \omterm{def:b1:gene-expression:ribosome}{ribossomos}.
\textbf{22.} $\num{3e9}\times 4 = \num{1.2e10}$ \omterm{def:b1:water-small-molecules:atp}{ATP} por dia
$= \qty{2e-14}{mol}$, \qty{1e-9}{J} por dia, ou seja, \qty{1.2e-14}{W}
--- \qty{6}{fW} por micrômetro cúbico.
\textbf{23.} $10^9$ \omterm{def:b1:water-small-molecules:atp}{ATP} por segundo são $\num{8.6e13}$ por dia: a
síntese proteica é um centésimo disso nessa \omterm{def:b1:cell-unit-of-life:cell}{célula} de renovação lenta
(numa bactéria em divisão ela é a metade).
\textbf{24.} Nenhum mensageiro maduro: os transcritos primários se
acumulam no núcleo e a \omterm{def:b1:proteins:peptide}{proteína} desaparece à medida que os mensageiros
dela se degradam, em poucas horas. A \omterm{def:b1:proteins:peptide}{proteína} bacteriana, cujo \omterm{def:b1:genomes:gene}{gene} não
tem \omterm{def:b1:genomes:gene}{íntrons}, não é afetada.
\textbf{25.} Cerca de 2100 \omterm{def:b1:water-small-molecules:atp}{ATP} por molécula --- 2000 para a tradução, e
cem para a parcela dela no transcrito; a primeira \omterm{def:b1:proteins:peptide}{proteína} depois de
\qty{1000}{s} de transcrição, um minuto de processamento e exportação, e
\qty{125}{s} de tradução: cerca de vinte minutos.
\end{solution}
